\subsection{Selected 1D Forward Model}
\label{subsec:selected-1d-forward-model}
\label{subsubsec:cross-sectional-reduction-selected-1d-model}

The model hierarchy identifies the 1D tier as the selected reduced forward map
for the numerical study. This subsection now records that selected model after
the hierarchy has established the 3D, multiscale, 2D, 1D, and 0D roles. The 1D
model reduces the spatial dimension by assuming that the vessel has a
distinguished axial direction and that pressure and velocity can be averaged
over cross-sections. This keeps pulse-wave and pressure-drop structure while
discarding secondary flow, separation, recirculation, and resolved wall traction.
The reduction follows standard compliant-artery and dimension-reduced modeling
arguments
\cite[Sec.~2, pp.~253--257]{FormaggiaEtAl2003OneDimensionalBloodFlow}
\cite[Ch.~3, Secs.~3.2 and~3.6]{KopplHelmig2023DimensionReducedModeling}
\cite[Sec.~3, pp.~24--26]{ShiEtAl2011BloodFlowModelReview}.

\subsubsection{Cross-Sectional Variables}

The reduced variables $A$, $Q$, $\bar u$, and $p_{\mathrm g}$ should be read as
cross-sectional outputs of the continuum fields under additional averaging,
geometry, wall-law, and profile assumptions, not as independent replacements for
the blood-dynamics setting.

\begin{figure}[!htb]
    \centering
    \captionsetup{aboveskip=2pt,belowskip=2pt}
    \figtikz{compliant-vessel}
    \caption{Compliant-vessel notation for the 1D reduction. The current
    configuration $\OmgBt$ is represented by a straightened averaging domain:
    section averages over $S(z,t)$ define $A(z,t)$, $Q(z,t)$, $\bar u(z,t)$,
    and $p_{\mathrm g}(z,t)$, while $R(z,t)$ records the radius entering the
    wall and geometry closures. The figure is a bookkeeping reduction from
    continuum fields to retained 1D variables, not a resolved wall-traction or
    secondary-flow model.}
    \label{fig:compliant-vessel-reduction}
\end{figure}

\begin{assumption}[Straight single-vessel geometry]
    \label{ass:1d-straight-vessel-geometry}
    The selected 1D model assumes a vessel of length $L$ aligned with the
    $z$-axis and a circular cross-section with radius $R(z,t)$. For
    $0<z<L$, $S(z,t)$ is the interior cross-section used for averaging, not
    a boundary component of the open lumen.
\end{assumption}

\begin{definition}[Section-averaged 1D fields]
    \label{def:section-averaged-1d-fields}
    The retained reduced variables are
    \begin{flalign*}
        \quad A(z,t) &\defined \int_{S(z,t)} 1 \, dS, &&
            \text{cross-sectional area}, \\
        \quad Q(z,t) &\defined \int_{S(z,t)} u_z \, dS, &&
            \text{volumetric flow rate}, \\
        \quad \bar u(z,t) &\defined \frac{Q(z,t)}{A(z,t)}, &&
            \text{mean axial velocity}, \\
        \quad p_{\mathrm g}(z,t) &\defined \frac{1}{A(z,t)}
            \int_{S(z,t)} p_{3D,\mathrm g} \, dS, &&
            \text{mean gauge pressure}.
    \end{flalign*}
    Here $p_{3D,\mathrm g}$ is continuum pressure in the wall-law gauge
    frame. The pressure-reference convention for reduced ratios is fixed
    separately in Convention~\ref{conv:appendix-pressure-reference-ratio}.
\end{definition}

\subsubsection{Closed Area--Flow Model and Wall Law}

\begin{definition}[Classical no-slip 1D baseline]
    \label{def:classical-no-slip-1d-baseline}
    The classical no-slip 1D baseline used as a model-family member starts from
    the fixed-wall continuum trace $\vu=0$ on the wall boundary
    $\Gamma_w\subset\pOmgB$, together with separate inlet and outlet data on
    $\Gamma_{\mathrm{in}}$ and $\Gamma_{\mathrm{out}}$. After cross-sectional
    averaging, the no-slip antecedent is represented by the parabolic
    profile/friction convention $\alpha_{\mathcal P}=4/3$ and
    $g_{\mathcal P}=4$. In the implementation manifests this member is recorded
    as \texttt{classical-1d-no-slip}; it keeps the same reduced state variables
    and selected Canic--Koiter pressure--area wall law as the extended 1D
    framework, but omits the Canic effective-alpha variable-radius correction.
    The descriptor \texttt{canic-extended-1d} records the current
    extended-model member.
\end{definition}

\begin{definition}[Closed 1D area-flow model]
    \label{def:closed-1d-area-flow-model}
    For $\mathcal W=\mathcal W_{\mathrm{elastic1D}}$ and an incompressible
    fluid with density $\rho$, the working compliant single-vessel model is
    written with the composite axial pressure derivative
    \begin{flalign*}
        \quad
        \frac{d}{dz}p_{\mathrm g}(A(z,t),z,t)
        =
        \pd_A p_{\mathrm g}(A,z,t) \pd_z A
        +
        \left.\pd_z p_{\mathrm g}(A,z,t)\right|_A. &&
    \end{flalign*}
    The reduced balance law is
    \begin{flalign*}
        \quad \pd_t A
        + \pd_z Q
        &= 0, \numberthis\label{eq:1d-mass-aq} && \\
        \quad \pd_t Q
        + \pd_z
            \left(\alpha \frac{Q^2}{A}\right)
        + \frac{A}{\rho}\frac{d}{dz}p_{\mathrm g}(A(z,t),z,t)
        &= - C_f(z,A,Q) \frac{Q}{A}.
        \numberthis\label{eq:1d-momentum-aq} &&
    \end{flalign*}
    The parameter $\alpha$ is the momentum-flux correction factor induced by
    the selected cross-sectional velocity profile, $C_f$ is the wall-friction
    closure, and $p_{\mathrm g}(A,z,t)$ is the wall-law gauge pressure
    closing the reduced system. This is the reduced
    $\mathcal W_{\mathrm{elastic1D}}$ member of the wall-model axis: wall
    response is represented through the pressure--area law rather than through a
    separate structural PDE
    \cite[Sec.~2, pp.~253--255]{FormaggiaEtAl2003OneDimensionalBloodFlow}
    \cite[Ch.~3, Sec.~3.2, pp.~40--47]{KopplHelmig2023DimensionReducedModeling}.
\end{definition}

\begin{assumption}[Velocity-profile and rheology closure]
    \label{ass:1d-first-run-alpha-friction}
    A closed 1D area-flow model must declare the momentum-flux correction
    factor $\alpha$, the wall-friction closure $C_f$, any rheology descriptor
    $\mathcal R$, and the pressure--area wall closure $\mathcal W$. These
    choices determine how cross-sectional profile effects, wall friction,
    viscosity or effective viscosity, and wall compliance enter the reduced
    equations. Numerical-record details such as solver coordinates,
    shear-rate proxies, and discrete source terms are deferred to
    Appendix~\ref{app:numerical-methods-details}.
\end{assumption}

\begin{definition}[Selected Canic--Koiter pressure--area wall law]
    \label{def:1d-elastic-wall-law}
    The selected elastic wall closure is the Canic--Koiter thin-wall
    pressure--area law. In conventional area notation it may be written
    \begin{flalign*}
        \quad p_{\mathrm g}(A,z,t)
        =
        p_{\mathrm{ext}}(z,t)
        + \frac{\beta(z)}{A_0(z)}
            \left(\sqrt{A(z,t)}-\sqrt{A_0(z)}\right).
        \numberthis\label{eq:1d-wall-law} &&
    \end{flalign*}
    Here $A_0(z)$ is the reference area, $p_{\mathrm{ext}}$ is the external
    pressure, and $\beta(z)$ is a wall-stiffness parameter. Some
    implementations absorb $A_0^{-1}$ into $\beta$; the essential closure
    is that $p_{\mathrm g}$ is known once $A_0$, $\beta$, and
    $p_{\mathrm{ext}}$ are supplied. In the radius-squared coordinate
    $a=R^2$ used by the implemented solver, the selected member is
    \begin{flalign*}
        \quad
        p_{1,\mathrm g}(a,z,t)
        =
        p_{\mathrm{ext}}(z,t)
        +
        \frac{K}{R_0(z)^2}
        \left(\sqrt a-R_0(z)\right),
        \qquad
        K=\frac{Eh}{1-\sigma^2}. &&
    \end{flalign*}
    The diagnostic extended pressure then adds the Canic variable-radius
    correction
    \begin{flalign*}
        \quad
        p_{2,\mathrm g}(a,Q,z,t)
        =
        g_{\mathcal P}\rho\nu_{\mathrm{eff}}^{\mathcal R}
        \frac{Q}{a}\frac{R_0'(z)}{R_0(z)}. &&
    \end{flalign*}
    This $p_2$ term is a variable-radius viscous correction, not a replacement
    wall law. Resolved FSI would replace the algebraic pressure--area closure
    with structural unknowns and fluid--structure interface conditions.
\end{definition}

\subsubsection{Implemented Coordinates and Stenosis Geometry}

\begin{definition}[Physical-to-solver map for the implemented 1D state]
    \label{def:physical-to-solver-map}
    The continuum 1D model uses physical area $A_{\mathrm{phys}}$ and
    physical flow $Q_{\mathrm{phys}}$. The implemented radius-squared solver
    instead stores
    \begin{flalign*}
        \quad
        a=R^2,\qquad q=\frac{Q_{\mathrm{phys}}}{\pi}.
        &&
    \end{flalign*}
    Thus
    \begin{flalign*}
        \quad
        A_{\mathrm{phys}}=\pi a,\qquad
        Q_{\mathrm{phys}}=\pi q,\qquad
        \bar u=\frac{Q_{\mathrm{phys}}}{A_{\mathrm{phys}}}=\frac{q}{a}.
        &&
    \end{flalign*}
    With $K=Eh/(1-\sigma^2)$ and the Canic conservative-form convention
    $R_0^*=R_{\max}$, the selected wall law gives the solver elastic flux
    contribution recorded in Appendix~\ref{app:num-current-solver-operator}:
    \begin{flalign*}
        \quad
        \Psi_h(a)=\frac{K}{3\rho R_{\max}^2}a^{3/2},
        \qquad
        \pd_a\Psi_h(a)=\frac{K}{2\rho R_{\max}^2}\sqrt a.
        &&
    \end{flalign*}
    This map is the convention used by the Section~\ref{sec:resolved-3d-comparison}
    velocity and flow comparisons.
\end{definition}

\begin{definition}[Stenosis reference geometry]
    \label{def:1d-stenosis-reference-geometry}
    A stenosis is encoded through the reference radius or reference area. Let
    $R_{\mathrm{base}}(z)$ be the nominal healthy radius and $s(z)\in[0,1)$
    a fractional radius-reduction profile. Then
    \begin{flalign*}
        \quad
        R_0(z) &= R_{\mathrm{base}}(z)\big(1-s(z)\big),
        &
        A_0(z) &= \pi R_0(z)^2.
        \numberthis\label{eq:1d-stenosis-area} &&
    \end{flalign*}
    The baseline idealized stenotic-vessel profile is the smooth asymmetric
    kernel
    \begin{flalign*}
        \quad
        g(z) &= z-z_s+\kappa
        \exp\left(-\frac{(z-z_a)^2}{2\sigma_a^2}\right),
        &
        \psi(z) &= \exp\left(-\lambda g(z)^4\right),
        \numberthis\label{eq:1d-stenosis-kernel} && \\
        \quad
        s(z) &= s_{\max}\psi(z), \qquad 0\leq s_{\max}<1.
        \numberthis\label{eq:1d-stenosis-profile} &&
    \end{flalign*}
    In the default centimeter-scale simulation record, $z_a=2.5\,\mathrm{cm}$,
    $z_s=3.4\,\mathrm{cm}$, $\sigma_a=1\,\mathrm{cm}$,
    $\kappa=0.95\,\mathrm{cm}$, and
    $\lambda=50\,\mathrm{cm}^{-4}$.
    This geometry is the baseline idealized simulation geometry in this report,
    not a patient-specific plaque morphology or clinical calibration.
\end{definition}

\begin{remark}[Stenosis-profile and variable-radius boundary]
    \label{rmk:1d-stenosis-profile-scope}
    The profile in Eq.~\ref{eq:1d-stenosis-profile} is $C^\infty$ and
    rapidly localized rather than compactly supported. This smooth idealized
    geometry is the baseline for the reduced simulations in this report and is
    also useful for fixed-wall 3D studies: the same analytic radius profile
    supplies controlled derivatives, wall normals, severity variation, and
    repeatable mesh-refinement families without introducing segmentation noise
    or patient-specific morphology. Canic et al.\ show that standard 1D
    variable-geometry models may miss stenotic pressure and velocity trends
    unless additional variable-radius terms are included
    \cite[Abstract; Secs.~1--2]{CanicEtAl2024Extended1DStenosis}. Those
    corrections are included only by the \texttt{canic-extended-1d}
    forward-model descriptor; \texttt{classical-1d-no-slip} retains the same
    selected wall law but omits the Canic variable-radius corrections.
\end{remark}

\subsubsection{Conservative Form and Boundary Contract}

\begin{definition}[Conservative flux/source notation for the reduced wall law]
    \label{def:1d-conservative-flux-source-notation}
    For a reduced wall closure $\mathcal W$ admitting an elastic potential,
    define
    \begin{flalign*}
        \quad
        \mathbf U
        &=
        (A,Q)^\top, &&
        \mathbf F_{\mathcal W}(\mathbf U,z,t)
        =
        \begin{pmatrix}
            Q\\[2pt]
            \alpha Q^2/A+\Psi_{\mathcal W}(A,z,t)
        \end{pmatrix}, &&\\
        \quad
        \mathbf S_{\mathcal W,\mathcal R}(\mathbf U,z,t)
        &=
        \begin{pmatrix}
            0\\[2pt]
            S_{\mathrm{geom}}^{\mathcal W}(A,z,t)
            +S_f^{\mathcal R}(A,Q,z)
        \end{pmatrix}. &&
    \end{flalign*}
    The potential and sources are tied to the selected closure descriptors by
    \begin{flalign*}
        \quad
        \pd_A \Psi_{\mathcal W}(A,z,t)
        &=
        \frac{A}{\rho} \pd_A p_{\mathrm g}^{\mathcal W}(A,z,t), &&\\
        \quad
        S_{\mathrm{geom}}^{\mathcal W}(A,z,t)
        &=
        \left.\pd_z \Psi_{\mathcal W}(A,z,t)\right|_A
        -
        \frac{A}{\rho}
        \left.\pd_z p_{\mathrm g}^{\mathcal W}(A,z,t)\right|_A,
        &&\\
        \quad
        S_f^{\mathcal R}(A,Q,z)
        &=
        -C_f^{\mathcal R}(z,A,Q)\frac{Q}{A}. &&
    \end{flalign*}
    For $\mathcal W_{\mathrm{elastic1D}}$,
    $p_{\mathrm g}^{\mathcal W}$ is the wall law in
    Definition~\ref{def:1d-elastic-wall-law}. In component form,
    \begin{flalign*}
        \quad \pd_t A + \pd_z Q &= 0,
        \numberthis\label{eq:1d-conservative-mass} && \\
        \quad \pd_t Q
        + \pd_z \left(
            \alpha\frac{Q^2}{A}+\Psi_{\mathcal W}(A,z,t)
        \right)
        &=
        S_{\mathrm{geom}}^{\mathcal W}(A,z,t)
        +S_f^{\mathcal R}(A,Q,z),
        \numberthis\label{eq:1d-conservative-momentum} &&
    \end{flalign*}
    The wall descriptor $\mathcal W$ determines the pressure--area law,
    elastic potential, wave speed, and geometry source; the rheology descriptor
    $\mathcal R$ determines the viscosity or friction closure. A resolved FSI
    model is therefore not obtained by merely relabeling
    $S_{\mathrm{geom}}$: it introduces structural state and interface laws.
    The displayed split records the balance-law objects needed by finite-volume
    schemes, while a computed realization may still use a particular
    well-balanced or source-split discretization
    \cite{LeVeque2002FiniteVolumeMethodsForHyperbolicProblems}
    \cite[Summary; Secs.~2.1--2.3; Secs.~5--6]{BoileauEtAl2015Benchmark}
    \cite{FormaggiaEtAl2007FSICoupling}.
\end{definition}

\begin{definition}[Single-vessel boundary data]
    \label{def:1d-single-vessel-boundary-data}
    For a single-vessel solve on $z\in[0,L]$, prescribe one inlet condition
    and one outlet condition:
    \begin{flalign*}
        \quad Q(0,t) &= Q_{\mathrm{in}}(t)
        \quad \text{or} \quad
        \bar u(0,t) = \bar u_{\mathrm{in}}(t), && \\
        \quad p_{\mathrm g}(L,t) &= p_{\mathrm{out},\mathrm g}(t)
        \quad \text{or} \quad
        p_{\mathrm g}(L,t) = p_{\mathrm{WK,g}}(Q(L,\cdot),t). &&
    \end{flalign*}
    The second outlet option denotes a terminal Windkessel-style relation. The
    numerical study uses prescribed $Q_{\mathrm{in}}(t)$ and
    $p_{\mathrm{out},\mathrm g}(t)$ because they give an unambiguous reduced
    single-vessel contract.
\end{definition}

\begin{definition}[Characteristic speeds for the reduced elastic subsystem]
    \label{def:1d-characteristic-speeds}
    Let $\lambda_\pm$ denote the two characteristic speeds of the
    positive-area elastic 1D subsystem obtained from the homogeneous
    $(A,Q)$ balance law. Their formula and assumptions are given in
    Appendix~\ref{app:num-elastic-potential-form}.
\end{definition}

\begin{assumption}[Subcritical boundary regime]
    \label{ass:1d-subcritical-boundary-regime}
    The baseline boundary treatment is used only in the subcritical regime
    $\lambda_-<0<\lambda_+$, where one characteristic enters through the inlet
    and one characteristic enters through the outlet.
\end{assumption}

Under this regime, a finite-volume scheme constrains the incoming information at
each boundary and leaves the outgoing state determined by the interior solve.
This is the reduced-model analogue of declaring enough boundary data for the
selected equation class; it is not, by itself, a validation of a particular
boundary discretization.
